\chapter{把算法接到 ROS 2 驱动器和实时循环}
\label{ch:software}

控制公式只规定理想的输入输出关系，工程实现还必须规定数据属于哪个时刻、哪个关节、哪个坐标系，以及谁拥有写命令的权利。本章以固定基座机械臂为例，把数学变量和软件接口逐个对应起来。软件接口说明按 ROS 2 Humble 体系组织；不同发行版和驱动器的实现细节不能直接混用。

\section{先定义接口 再选择框架}

一个可复用控制器至少需要三类结构：测量、参考和输出。测量是 $\{q,\dot q,\tau_{\rm meas},t_{\rm sample},\mathrm{valid}\}$；参考是 $\{q_d,\dot q_d,\ddot q_d,T_d,\mathrm{mode}\}$ 中本控制器实际需要的量；输出是明确单位的关节命令。不是每个机器人都能提供真实关节力矩，也不是每个控制器都需要全部字段。

\begin{table}[htbp]\centering\small
\begin{tabularx}{\linewidth}{p{28mm}XX}\toprule
边界 & 必须写清楚的约定 & 数学意义\\\midrule
编码器到状态 & 零位 方向 角度单位 时间戳 & $q$ 不是未经转换的计数值\\
动力学到力矩 & 电机侧或关节侧 重力方向 & $\tau$ 与 $M\ddot q$ 同单位\\
任务到关节 & 末端参考点 坐标轴 行顺序 & $\dot x=J\dot q$ 必须同约定\\
控制器到驱动 & 接口模式 内置增益 补偿所有权 & 避免同一 PD 或重力补偿做两次\\
线程之间 & 快照一致性 数据年龄 & 防止把不同拍关节拼成一个状态\\\bottomrule
\end{tabularx}
\caption{这些约定是控制算法的一部分，不能只写在口头说明里。}
\end{table}

例如关节消息顺序为 $[j_3,j_1,j_2]$，值为 $[0.3,0.1,0.2]$，模型顺序为 $[j_1,j_2,j_3]$，正确输入应是 $[0.1,0.2,0.3]$。直接把消息数组复制到模型向量不会报矩阵维度错误，却会让正运动学和重力补偿全部算错。

\section{ROS 2 的节点 消息和坐标树}

\textbf{节点（node）}是逻辑计算单元，不等于操作系统进程；多个节点可以组合在同一进程中。\textbf{话题（topic）}提供发布订阅流；\textbf{服务（service）}适合请求与响应；\textbf{动作（action）}适合有持续时间、进度反馈和取消语义的任务，例如执行一条轨迹。机器人模型与坐标变换不是实时运动命令。

\textbf{URDF}描述连杆、关节、几何、惯量和部分限位；\textbf{xacro}帮助复用模型描述；\textbf{tf2}维护带时间信息的坐标变换树。某个变换数值正确但时间不匹配，仍可能把移动相机的图像放到错误位置。仿真时相关节点应使用一致的时间基准；实际时钟和仿真时钟不能随意相减。

\textbf{QoS（quality of service，服务质量）}规定可靠性、历史队列和持久性等通信策略。可靠传输关注是否重传，实时控制还关注这份数据是否已经过期；队列越长不一定越好。订阅端所请求的 QoS 与发布端提供的 QoS 不兼容时，可能无法连接。可以检查端点信息和事件，而不是只看话题名字存在。

下列命令用于只读检查接口；其中话题名来自具体系统，不能当成所有 ROS 2 机器人的标准控制接口：
\begin{lstlisting}[language=bash]
ros2 node list
ros2 topic list -t
ros2 topic info /control/move_j --verbose
ros2 topic type /control/move_j
ros2 interface show sensor_msgs/msg/JointState
ros2 topic echo /feedback/joint_states --once
\end{lstlisting}
发布运动之前，要先查真实消息类型、关节名、订阅者和接口语义。消息叫 \texttt{JointState} 并不证明它只能用作状态，也不证明任意填一个 \texttt{position} 数组就能执行；有些厂家把标准消息用于私有命令约定，解释权在驱动实现。

\section{ros2\_control 的一个周期}

ros2\_control 把硬件通信与控制算法分离。硬件提供\textbf{状态接口}和\textbf{命令接口}；控制器读取状态并写命令；控制器管理器组织生命周期和周期更新。典型循环为读取硬件、计算控制、写出命令。官方 Humble 文档把 \texttt{update\_rate} 定义为这条更新循环的频率，同时强调需要控制抖动。\cite{ros2control}

硬件插件不只是两个孤立的函数。它还要完成初始化、资源分配、激活、停用、错误处理和命令模式切换。控制器激活时应把位置目标与当前状态对齐，并按控制律处理积分器和滤波器状态；仅把所有内部量清零不一定无扰。

\subsection{轨迹控制器到底把什么交给驱动器}

Humble 的关节轨迹控制器按照时间对航点插补。位置命令接口通常转发插补后的位置；速度或力矩接口可通过配置的 PID 把位置和速度跟踪误差映射为命令。可用接口组合和状态反馈要求应以对应版本文档为准。\cite{jointtrajectory}

因此，同样一条关节轨迹，在位置接口和力矩接口下的闭环结构不同。位置接口下面通常还有驱动器的内部伺服；力矩接口则要求上层承担更多动态控制职责。把上层输出的 $-K_pe-K_d\dot e$ 再送进另一个含同样 PD 的接口，可能叠加刚度和阻尼。

以常见组合式关节接口为例，若驱动器实现
\begin{equation}
\tau=K_{p,\rm drv}(q_r-q)+K_{d,\rm drv}(v_r-\dot q)+\tau_{\rm ff},
\end{equation}
则上位机可以选择让驱动器承担局部 PD，自己计算前馈；也可以在协议允许且行为已确认时，使用合适设置把主要力矩律放在上位机。必须明确两个层级各自做什么。协议中的增益数值范围、编码比例、饱和和电流限幅，不能由这个公式反推出来。

\section{MoveIt 与动力学库的分工}

MoveIt 的模型、规划场景和规划管线负责运动规划相关工作；ros2\_control 组织执行接口；Pinocchio 这类动力学库负责计算运动学、惯量、非线性项和雅可比。库能很快计算公式，但不会自动决定你的控制器应该用哪一种误差、哪一种摩擦补偿或哪一种安全停机策略。

\begin{table}[htbp]\centering\small
\begin{tabularx}{\linewidth}{p{34mm}XX}\toprule
配置或数据 & 主要用途 & 与其它部分的边界\\\midrule
URDF 及碰撞几何 & 机构 轴向 惯量 模型限位 & CAD 惯量不等于已完成辨识\\
SRDF & 规划组 命名姿态 允许碰撞对 & 禁止滥用允许碰撞掩盖几何错误\\
运动学配置 & IK 插件 超时 搜索策略 & 七轴不意味着某个插件必然不能用\\
规划限值 & 速度 加速度等 & 与驱动器硬限值保持一致和余量\\
执行配置 & 控制器名字 动作接口 关节集合 & 规划成功不等于执行已成功\\
工具附件 & 相机或抓持物的碰撞形状 & 只标 TCP 不会自动添加碰撞体\\\bottomrule
\end{tabularx}
\caption{软件配置表达的是算法前提，出现问题时要回到其物理含义。}
\end{table}

选择 IK 插件时，应按实际模型、限位、目标分布、初值和超时，比较成功率及耗时。原稿“七轴不要用 KDL”的说法过于绝对；适用性应由具体任务与求解表现判断。

\subsection{模型配置维度和速度维度}

配置向量的维度记为 $n_q$，切空间速度维度记为 $n_v$。普通有界转动关节各占一个配置量和一个速度量；若某库用正余弦表示无限转动关节，配置量可能有两个。浮动基座常用位置加单位四元数存七个数，速度则是三维线速度加三维角速度，共六个数。

因此一般的配置更新是 $q_{k+1}=\operatorname{integrate}(q_k,v_kT_s)$，不是不加区分地执行 $q+vT_s$。对本书固定基座、普通角度坐标的 2R 模型，普通加法可以成立；推广时需重新确认库的表示约定。

同样，\texttt{LOCAL}、\texttt{WORLD} 和\texttt{LOCAL\_WORLD\_ALIGNED} 的区别不仅是坐标轴旋转，还涉及空间速度参考点。控制末端点的线速度时，需要末端参考点的雅可比；绕世界原点表达的空间旋量，其前三项未必等于末端点的实际线速度。第3章的旋量变换给出了这一区别的数学原因。

\section{实时性用时间预算解释}

设控制周期为 $T_s=5\,\mathrm{ms}$，对应 $200\,\mathrm{Hz}$。一个教学预算为读取 $0.7\,\mathrm{ms}$、计算 $0.9\,\mathrm{ms}$、写出 $0.5\,\mathrm{ms}$、其它固定开销 $0.2\,\mathrm{ms}$，总计 $2.3\,\mathrm{ms}$，余量为 $2.7\,\mathrm{ms}$。若偶发一次文件写入耗时 $8\,\mathrm{ms}$，平均耗时依旧可能很好看，但这一拍已经错过截止时间。

\begin{figure}[htbp]\centering
\includegraphics[width=0.93\linewidth]{figures/rt_timing.pdf}
\caption{五毫秒周期的时间预算示例。绝对时间调度可以避免把每次计算时间累积进下一次周期，但不能消除过载、抢占和通信阻塞。}
\end{figure}

\textbf{周期}是理想更新间隔；\textbf{唤醒抖动}是实际开始时刻偏离计划开始时刻的量；\textbf{执行耗时}是本轮代码实际用了多久；\textbf{端到端延迟}从传感器采样一直算到执行器产生作用。它们相关但不相同。

若每轮先计算 $0.8\,\mathrm{ms}$ 再固定睡眠 $5\,\mathrm{ms}$，实际周期至少为 $5.8\,\mathrm{ms}$，频率约为 $172.4\,\mathrm{Hz}$。采用绝对时刻 $t_k=t_0+kT_s$ 作为下一次唤醒目标，可以避免这类漂移；某轮超时后则必须记录并按设计处理，不能用“快速补跑多次旧命令”伪装没超时。

热路径中应避免未经界定的动态分配、可能阻塞的输入输出和无上界等待。采用预分配存储、一致的数据快照、有界队列和经过验证的线程交互。无锁本身也不是正确性的保证：内存顺序和读写一致性仍需证明。普通 ROS 2 话题并不自动提供硬实时保证，但能否满足具体周期取决于体系结构与测量，不能用“所有 ROS 消息永远不能进入实时系统”替代分析。

\section{CAN 带宽先算账再谈频率}

只做一个明确假设的估算：每个关节每周期一帧命令、一帧反馈，每帧连同开销平均按 $130\,\mathrm{bit}$ 估计，总线速率 $1\,\mathrm{Mbit/s}$。七关节、$200\,\mathrm{Hz}$ 时
\begin{equation}
B=7\times2\times200\times130=364000\,\mathrm{bit/s},
\end{equation}
平均占用约 $36.4\%$。改为 $1000\,\mathrm{Hz}$ 时，需求为 $1.82\,\mathrm{Mbit/s}$，已超过单路总线能力。这个算例不代表某款机械臂的实际协议：位填充、帧格式、打包、多路总线、额外状态报文和错误重发都会改变结果。

平均带宽还不等于最坏延迟。仲裁会让低优先级帧等待；同一拍七个关节反馈不一定同时到达。需要记录每轴反馈年龄，必要时做有说明的状态外推，并给异常数据设置明确的处理行为。总线速率高，不代表控制闭环带宽也相同。

\section{编码器 电气接口和力矩单位}

编码器是传感器，不会自己驱动机械臂。增量式编码器的 A、B 两路有相位差，用边沿先后判断方向；Z 通常是每转一次的参考标记。若每路每转 $N=2500$ 个周期，采用四倍频计数，一转有 $10000$ 个计数，对应
\begin{equation}
\Delta q=2\pi/10000=0.00062832\,\mathrm{rad}=0.036^\circ.
\end{equation}
接在电机端时，还需按减速比换算到关节端，并考虑零位与传动弹性。标称分辨率不等于实际准确度。

A+/A-、B+/B-、Z+/Z- 通常表示差分信号对。DB9 只是连接器形状，把导线焊到 DB9 上不会自动变成 RS232 串行数据。差分接收、正交计数、电源和逻辑电平匹配是不同功能；接口类型必须根据设备输出电路确定。普通串口适配器不会仅凭插头形状就读出角度。

力矩换算也要保留层级。电机轴 $\tau_m=K_ti_q$，理想减速器输出约为 $\tau_j=\eta N\tau_m$。如果驱动反馈已经是关节侧力矩，再乘一次 $N$ 就会放大几十或上百倍。厂家的 $k,b,c$ 等系数名称没有跨厂家的统一含义，要把每个系数的单位、标定依据和已经包含的变换写下来。

\section{控制状态机与命令所有权}

一个教学状态机可以包含未激活、保持、轨迹、拖动、受控停止和故障。状态转移同时改变参考生成器、控制律和内部状态，不能只把字符串变量换掉。每个关节在同一时刻只能有明确的一套最终命令来源；“规划器还在发位置，同时拖动程序又发 MIT 力矩”不是控制融合，而是竞争写入。

停止也不是把所有数值写成零。重力存在时，零力矩可能导致下落；位置置零可能发起大幅回零运动。正常减速停止、失联处理、硬件紧急停止和断电保持属于不同机制，其具体动作由机械结构、抱闸和驱动器设计决定。教材不规定某款设备通用的使能延时或抱闸时序。

\section{计算练习与解答}

\subsection*{题一 总线预算}
采用上述每关节收发各一帧、每帧 $130\,\mathrm{bit}$ 假设，七关节运行 $500\,\mathrm{Hz}$，占用多少带宽？

解：$7\times2\times500\times130=910000\,\mathrm{bit/s}$，即 $91\%$。这尚未包含额外报文和重发，不能只因为小于 $100\%$ 就认定有充分余量。

\subsection*{题二 电机侧与关节侧}
已知 $K_t=0.08\,\mathrm{N\,m/A},i_q=3\,\mathrm A,N=50,\eta=0.9$。求电机和关节侧力矩。

解：$\tau_m=0.08\times3=0.24\,\mathrm{N\,m}$，$\tau_j=0.9\times50\times0.24=10.8\,\mathrm{N\,m}$。这个简化静态换算不含方向相关摩擦和动态损耗。

\subsection*{题三 周期与数据年龄}
控制周期为 $5\,\mathrm{ms}$，最新反馈采样时间距当前已过去 $12\,\mathrm{ms}$。若关节可能以 $0.5\,\mathrm{rad/s}$ 运动，忽略加速度，位置可能相差多少？

解：$0.5\times0.012=0.006\,\mathrm{rad}\approx0.344^\circ$。它跨越了 $2.4$ 个周期，不能只检查“这一拍是否收到一条消息”，还要检查消息对应的真实采样时刻。

\section*{本章与原稿的对应}
本章整合 [A，第67--74页] 与 [B，第100--111、116--120页]，并补充编码器接口与带宽计算。纠正节点等于进程、ROS 通信天然实时、七轴必须排除某个 IK 插件、停机等于写零以及不同 MoveIt 版本接口混用等问题。
